%% ============================================================
%% cap04.tex — Capítulo 4: Implementación en Java
%% Programación Orientada a Objetos: Diseño e Implementación
%% de un Sistema de Información Académico-Administrativo
%% para la Universidad Distrital Francisco José de Caldas
%%
%% Autor: Roberto Albeiro Pava-Díaz, Ph.D.
%% ============================================================

\chapter{Implementación en \java}
\label{cap:java}


\noindent
El tránsito del diagrama UML a la línea de código representa, quizás, el momento más
revelador del proceso de aprendizaje en \poo. Es allí donde las abstracciones del
modelado se confrontan con las restricciones concretas de un lenguaje, de un compilador
y, sobre todo, de un dominio que nunca es tan ordenado como el papel sugiere. El dominio
del estatuto de la \ud\ no constituye excepción: la jerarquía de unidades académicas,
la multiplicidad de vínculos docentes y la granularidad del sistema de evaluación
formativa imponen decisiones de diseño que no se resuelven en el pizarrón.

Este capítulo expone la implementación completa del sistema en \java, lenguaje que
ha acompañado la enseñanza de la orientación a objetos durante más de dos décadas y
cuya plataforma JVM ofrece un ecosistema maduro para aplicaciones de gestión institucional.
Se examinan la arquitectura por capas, las clases del dominio, los patrones de diseño
adoptados, el manejo de excepciones, la interfaz gráfica con JavaFX y la persistencia
mediante serialización JSON. Cada sección incluye código fuente completo con comentarios
Javadoc, de manera que el lector pueda trasladar los listados directamente al entorno
de desarrollo integrado.\footnote{El proyecto completo, estructurado como proyecto Maven
listo para Visual Studio Code, se encuentra disponible en el directorio
\texttt{proyecto\_java\_estatuto/} que acompaña este libro. Su compilación requiere
Java~17 o superior y Maven~3.8+.}

% ─────────────────────────────────────────────────────────────────────────────
\section{Arquitectura del proyecto \java}
\label{sec:java-arquitectura}
% ─────────────────────────────────────────────────────────────────────────────

\noindent
El principio de separación de responsabilidades (\textit{separation of concerns})
constituye el fundamento sobre el cual se erige la arquitectura del sistema.
Se adopta el patrón arquitectónico Modelo-Vista-Controlador (MVC), que distribuye
la lógica del sistema en tres capas bien delimitadas, reduciendo el acoplamiento y
facilitando la evolución independiente de cada componente.\footnote{La adopción de MVC
en aplicaciones de gestión académica universitaria ha sido documentada, entre otros,
por \textcite{Gamma1994} y, en el contexto latinoamericano, por trabajos de modelado
institucional que abordan la complejidad de los sistemas de información universitarios.}

\subsection{Estructura de paquetes}

La tabla~\ref{tab:paquetes-java} describe los cinco paquetes que conforman el entramado
operativo del sistema, su responsabilidad asignada y las clases representativas que
cada uno alberga.

\begingroup
\small
\begin{longtable}{@{}L{5cm} L{3.5cm} L{6cm}@{}}
	% --- Título y encabezado de la primera página ---
	\caption{Paquetes del sistema y sus responsabilidades.} \label{tab:paquetes-java} \\
	\toprule
	\textbf{Paquete} & \textbf{Capa MVC} & \textbf{Responsabilidad principal} \\
	\midrule
	\endfirsthead
	
	% --- Encabezado para las páginas siguientes ---
	\multicolumn{3}{c}{{\bfseries \tablename\ \thetable{} -- Continuación}} \\
	\toprule
	\textbf{Paquete} & \textbf{Capa MVC} & \textbf{Responsabilidad principal} \\
	\midrule
	\endhead
	
	% --- Pie de tabla para páginas intermedias ---
	\midrule
	\multicolumn{3}{r}{{Continúa en la siguiente página...}} \\
	\bottomrule
	\endfoot
	
	% --- Pie de tabla final ---
	\bottomrule
	\endlastfoot
	
	% --- Datos de la tabla ---
	\texttt{co.edu.udistrital.} \texttt{estatuto.modelo} &
	Modelo &
	Clases de dominio: entidades académico-administrativas, enumeraciones y objetos de valor \\
	\midrule
	
	\texttt{co.edu.udistrital.} \texttt{estatuto.controlador} &
	Controlador &
	Lógica de negocio, validaciones, orquestación de operaciones CRUD \\
	\midrule
	
	\texttt{co.edu.udistrital.} \texttt{estatuto.vista} &
	Vista &
	Componentes JavaFX: ventanas, formularios, tablas, menús \\
	\midrule
	
	\texttt{co.edu.udistrital.} \texttt{estatuto.persistencia} &
	Infraestructura &
	Acceso a datos mediante patrón DAO; serialización/deserialización JSON con Gson \\
	\midrule
	
	\texttt{co.edu.udistrital.} \texttt{estatuto.util} &
	Transversal &
	Validadores, formateadores de fecha, constantes del sistema \\
\end{longtable}
\endgroup

\noindent
Los paquetes \texttt{patron} y \texttt{excepcion} complementan la estructura anterior;
el primero aloja los patrones de diseño implementados (Factory Method, Strategy, Observer)
y el segundo, la jerarquía de excepciones del dominio. Esta decisión de separar los
patrones en un paquete propio ---discutible desde la perspectiva de la cohesión--- se
justifica didácticamente: facilita la identificación y el estudio aislado de cada patrón.

\subsection{JavaFX como \textit{framework} de interfaz gráfica}

La elección de JavaFX como \textit{framework} GUI obedece a tres razones concurrentes.
Primera, su integración nativa con el ecosistema \java\ permite que los estudiantes
reutilicen conceptos ya adquiridos (herencia, eventos, hilos) sin abandonar el lenguaje.
Segunda, el modelo de componentes de JavaFX ---basado en un grafo de escena
(\textit{scene graph})--- refuerza la comprensión de las jerarquías de objetos:
cada nodo es un objeto que hereda de \texttt{javafx.scene.Node}, exhibiendo
polimorfismo en cada interacción con la interfaz. Tercera, la separación entre
la definición de la interfaz (FXML) y su lógica (controlador JavaFX) ilustra,
en un contexto práctico, el principio de separación de responsabilidades que se
defiende en este libro.\footnote{JavaFX fue desacoplado del JDK estándar a partir
de Java~11. Para este proyecto se utiliza JavaFX~21 LTS como dependencia Maven,
lo que garantiza compatibilidad con Java~17 y Java~21.}

% ─────────────────────────────────────────────────────────────────────────────
\section{Clases fundamentales del dominio}
\label{sec:java-dominio}
% ─────────────────────────────────────────────────────────────────────────────

\noindent
El modelado del dominio traduce las entidades identificadas en el capítulo~\ref{cap:modelado}
a clases \java\ concretas. La fidelidad a ese modelo no implica, sin embargo, una
correspondencia mecánica uno a uno: la implementación introduce decisiones técnicas
---constructores, validadores, sobreescritura de \texttt{equals}--- que el diagrama
UML soslaya deliberadamente. A continuación se presentan las clases agrupadas según
su función en el dominio.

\subsection{Clases de entidad organizacional}
\label{subsec:entidad-organizacional}

La clase abstracta \texttt{UnidadAcademica} constituye la raíz de la jerarquía
de entidades organizacionales. Toda vez que el \acuerdo\ define Facultades, Escuelas,
Centros e Institutos como unidades académico-administrativas con atributos comunes
---nombre, código y fecha de creación---, resulta natural encapsular esos atributos
en una superclase que los concrete y obligue a sus subclases a declarar su tipo mediante
el método abstracto \texttt{getTipo()}.

\begin{lstlisting}[style=javaStyle, caption={Clase abstracta \texttt{UnidadAcademica}}, label={lst:unidad-academica}]
package co.edu.udistrital.estatuto.modelo;

import java.time.LocalDate;
import java.util.Objects;

/**
 * Clase abstracta que representa una unidad académico-administrativa
 * de la Universidad Distrital Francisco José de Caldas.
 *
 * <p>Conforme al Acuerdo 004 de 2025 (Estatuto General), toda unidad
 * académica posee un nombre, un código institucional, una fecha de
 * creación oficial y un director designado por el órgano competente.</p>
 *
 * @author Roberto Albeiro Pava-Díaz
 * @version 1.0
 */
public abstract class UnidadAcademica {

    /** Nombre oficial de la unidad académica. */
    private String nombre;

    /** Código institucional único asignado por la Rectoría. */
    private String codigo;

    /** Fecha de creación o reconocimiento oficial de la unidad. */
    private LocalDate fechaCreacion;

    /** Nombre completo del director o directora vigente. */
    private String director;

    /**
     * Construye una unidad académica con los datos básicos.
     *
     * @param nombre        nombre oficial de la unidad
     * @param codigo        código institucional (no nulo, no vacío)
     * @param fechaCreacion fecha de creación oficial
     * @param director      nombre del director vigente
     * @throws IllegalArgumentException si el código es nulo o vacío
     */
    protected UnidadAcademica(String nombre, String codigo,
                               LocalDate fechaCreacion, String director) {
        if (codigo == null || codigo.isBlank()) {
            throw new IllegalArgumentException(
                "El código de la unidad académica no puede ser nulo ni vacío.");
        }
        this.nombre = Objects.requireNonNull(nombre, "El nombre no puede ser nulo.");
        this.codigo = codigo.trim().toUpperCase();
        this.fechaCreacion = fechaCreacion;
        this.director = director;
    }

    /**
     * Retorna el tipo de unidad académica (Facultad, Escuela, Centro, etc.).
     * Cada subclase debe implementar este método con su denominación específica.
     *
     * @return cadena descriptiva del tipo de unidad
     */
    public abstract String getTipo();

    // ── Getters y Setters ──────────────────────────────────────────────────

    /** @return nombre oficial de la unidad */
    public String getNombre() { return nombre; }

    /**
     * Actualiza el nombre oficial de la unidad.
     * @param nombre nuevo nombre (no nulo, no vacío)
     */
    public void setNombre(String nombre) {
        if (nombre == null || nombre.isBlank())
            throw new IllegalArgumentException("El nombre no puede ser vacío.");
        this.nombre = nombre.trim();
    }

    /** @return código institucional */
    public String getCodigo() { return codigo; }

    /**
     * Actualiza el código institucional. Se normaliza a mayúsculas.
     * @param codigo nuevo código
     */
    public void setCodigo(String codigo) {
        if (codigo == null || codigo.isBlank())
            throw new IllegalArgumentException("El código no puede ser vacío.");
        this.codigo = codigo.trim().toUpperCase();
    }

    /** @return fecha de creación oficial */
    public LocalDate getFechaCreacion() { return fechaCreacion; }

    /** @param fechaCreacion nueva fecha de creación */
    public void setFechaCreacion(LocalDate fechaCreacion) {
        this.fechaCreacion = fechaCreacion;
    }

    /** @return nombre del director vigente */
    public String getDirector() { return director; }

    /** @param director nombre del nuevo director */
    public void setDirector(String director) { this.director = director; }

    // ── equals, hashCode, toString ─────────────────────────────────────────

    @Override
    public boolean equals(Object o) {
        if (this == o) return true;
        if (!(o instanceof UnidadAcademica u)) return false;
        return codigo.equals(u.codigo);
    }

    @Override
    public int hashCode() { return Objects.hash(codigo); }

    @Override
    public String toString() {
        return String.format("%s[codigo=%s, nombre=%s, director=%s]",
                             getTipo(), codigo, nombre, director);
    }
}
\end{lstlisting}

\noindent
La clase \texttt{Facultad} extiende \texttt{UnidadAcademica} e incorpora los atributos
propios de su nivel jerárquico: el decano o decana y la lista de escuelas adscritas.
La agregación de escuelas se implementa mediante una lista mutable que el controlador
gestiona externamente, de manera que la clase de dominio no asuma responsabilidades
de persistencia.\footnote{Esta separación entre el objeto de dominio y su repositorio
corresponde al patrón \textit{Domain Model} documentado por \textcite{fowler2002patterns}.}

\begin{lstlisting}[style=javaStyle, caption={Clase \texttt{Facultad}}, label={lst:facultad}]
package co.edu.udistrital.estatuto.modelo;

import java.time.LocalDate;
import java.util.ArrayList;
import java.util.Collections;
import java.util.List;
import java.util.Objects;

/**
 * Representa una Facultad de la Universidad Distrital.
 *
 * <p>Las Facultades son las unidades académico-administrativas de primer
 * nivel (Art. 24, Acuerdo 004 de 2025). Cada Facultad está encabezada
 * por un Decano o Decana y agrupa un conjunto de Escuelas.</p>
 *
 * @author Roberto Albeiro Pava-Díaz
 */
public class Facultad extends UnidadAcademica {

    /** Nombre completo del decano o decana vigente. */
    private String decano;

    /** Lista de escuelas adscritas a esta Facultad. */
    private final List<Escuela> escuelas;

    /**
     * Construye una Facultad con los datos mínimos requeridos.
     *
     * @param nombre        nombre oficial de la facultad
     * @param codigo        código institucional
     * @param fechaCreacion fecha de creación
     * @param director      director académico (equivalente al decano aquí)
     * @param decano        nombre del decano o decana
     */
    public Facultad(String nombre, String codigo,
                    LocalDate fechaCreacion, String director, String decano) {
        super(nombre, codigo, fechaCreacion, director);
        this.decano = Objects.requireNonNull(decano, "El decano no puede ser nulo.");
        this.escuelas = new ArrayList<>();
    }

    @Override
    public String getTipo() { return "Facultad"; }

    /**
     * Agrega una Escuela a la lista de escuelas de la Facultad.
     * No se permiten escuelas duplicadas (comparación por código).
     *
     * @param escuela escuela a agregar
     * @return {@code true} si se agregó, {@code false} si ya existía
     */
    public boolean agregarEscuela(Escuela escuela) {
        Objects.requireNonNull(escuela, "La escuela no puede ser nula.");
        if (escuelas.contains(escuela)) return false;
        return escuelas.add(escuela);
    }

    /**
     * Elimina una Escuela de la lista.
     *
     * @param codigo código de la escuela a eliminar
     * @return {@code true} si se eliminó
     */
    public boolean eliminarEscuela(String codigo) {
        return escuelas.removeIf(e -> e.getCodigo().equals(codigo));
    }

    /** @return vista no modificable de la lista de escuelas */
    public List<Escuela> getEscuelas() {
        return Collections.unmodifiableList(escuelas);
    }

    /** @return nombre del decano o decana */
    public String getDecano() { return decano; }

    /** @param decano nombre del nuevo decano */
    public void setDecano(String decano) { this.decano = decano; }

    @Override
    public String toString() {
        return String.format("Facultad[codigo=%s, nombre=%s, decano=%s, escuelas=%d]",
                             getCodigo(), getNombre(), decano, escuelas.size());
    }
}
\end{lstlisting}

\begin{lstlisting}[style=javaStyle, caption={Clases \texttt{Escuela}, \texttt{Centro} e \texttt{Instituto}}, label={lst:escuela-centro}]
package co.edu.udistrital.estatuto.modelo;

import java.time.LocalDate;
import java.util.*;

/**
 * Representa una Escuela adscrita a una Facultad.
 * Todos los docentes de la universidad pertenecen a una Escuela
 * (Art. 28, Acuerdo 004 de 2025).
 */
public class Escuela extends UnidadAcademica {

    private final List<Docente>  docentes;
    private final List<Caba>     cabas;

    public Escuela(String nombre, String codigo,
                   LocalDate fechaCreacion, String director) {
        super(nombre, codigo, fechaCreacion, director);
        this.docentes = new ArrayList<>();
        this.cabas    = new ArrayList<>();
    }

    @Override public String getTipo() { return "Escuela"; }

    public boolean agregarDocente(Docente d) {
        Objects.requireNonNull(d);
        return !docentes.contains(d) && docentes.add(d);
    }

    public boolean agregarCaba(Caba c) {
        Objects.requireNonNull(c);
        return !cabas.contains(c) && cabas.add(c);
    }

    public List<Docente> getDocentes() {
        return Collections.unmodifiableList(docentes);
    }

    public List<Caba> getCabas() {
        return Collections.unmodifiableList(cabas);
    }
}

// ─────────────────────────────────────────────────────────────────────────────

/**
 * Representa un Centro universitario con consejo y dirección propios
 * (Art. 30, Acuerdo 004 de 2025).
 */
class Centro extends UnidadAcademica {
    public Centro(String nombre, String codigo,
                  LocalDate fechaCreacion, String director) {
        super(nombre, codigo, fechaCreacion, director);
    }
    @Override public String getTipo() { return "Centro"; }
}

// ─────────────────────────────────────────────────────────────────────────────

/**
 * Representa un Instituto universitario con consejo y dirección propios
 * (Art. 32, Acuerdo 004 de 2025).
 */
class Instituto extends UnidadAcademica {
    public Instituto(String nombre, String codigo,
                     LocalDate fechaCreacion, String director) {
        super(nombre, codigo, fechaCreacion, director);
    }
    @Override public String getTipo() { return "Instituto"; }
}
\end{lstlisting}

\begin{lstlisting}[style=javaStyle, caption={Clase \texttt{Caba} (Comité Académico Básico de Área)}, label={lst:caba}]
package co.edu.udistrital.estatuto.modelo;

import java.util.Objects;

/**
 * Comité Académico Básico de Área (CABA).
 *
 * <p>Los CABAS articulan las funciones de docencia e investigación
 * al interior de cada Escuela, tal como lo dispone el Acuerdo 004 de 2025.
 * Cada CABA pertenece a un área de formación y es coordinado por un
 * docente designado por el Consejo de Escuela.</p>
 */
public class Caba {

    private String areaFormacion;
    private String coordinador;
    private String nombre;

    public Caba(String nombre, String areaFormacion, String coordinador) {
        this.nombre        = Objects.requireNonNull(nombre);
        this.areaFormacion = Objects.requireNonNull(areaFormacion);
        this.coordinador   = coordinador;
    }

    public String getNombre()        { return nombre; }
    public String getAreaFormacion() { return areaFormacion; }
    public String getCoordinador()   { return coordinador; }

    public void setAreaFormacion(String af) { this.areaFormacion = af; }
    public void setCoordinador(String c)    { this.coordinador   = c; }

    @Override
    public String toString() {
        return String.format("CABA[nombre=%s, area=%s, coordinador=%s]",
                             nombre, areaFormacion, coordinador);
    }
}
\end{lstlisting}

\subsection{Clases de actores}
\label{subsec:actores}

La comunidad universitaria, definida en el artículo~8 del \acuerdo, comprende docentes,
estudiantes, personal administrativo y egresados. Todos comparten atributos personales
básicos encapsulados en la clase abstracta \texttt{Persona}, cuya herencia evita la
duplicación de código y garantiza que todo actor del sistema sea identificable
de manera unívoca.

\begin{lstlisting}[style=javaStyle, caption={Clase abstracta \texttt{Persona} y enumeraciones de actor}, label={lst:persona}]
package co.edu.udistrital.estatuto.modelo;

import java.util.Objects;

/**
 * Clase abstracta que representa a cualquier persona vinculada
 * a la Universidad Distrital Francisco José de Caldas.
 *
 * @author Roberto Albeiro Pava-Díaz
 */
public abstract class Persona {

    private String nombre;
    private String identificacion;
    private String correo;

    protected Persona(String nombre, String identificacion, String correo) {
        this.nombre         = Objects.requireNonNull(nombre, "Nombre requerido.");
        this.identificacion = Objects.requireNonNull(identificacion, "ID requerido.");
        this.correo         = correo;
    }

    /** @return rol institucional de la persona (Docente, Estudiante, etc.) */
    public abstract String getRol();

    public String getNombre()         { return nombre; }
    public void   setNombre(String n) { this.nombre = n; }

    public String getIdentificacion()         { return identificacion; }
    public void   setIdentificacion(String i) { this.identificacion = i; }

    public String getCorreo()         { return correo; }
    public void   setCorreo(String c) { this.correo = c; }

    @Override
    public boolean equals(Object o) {
        if (this == o) return true;
        if (!(o instanceof Persona p)) return false;
        return identificacion.equals(p.identificacion);
    }

    @Override public int hashCode() { return Objects.hash(identificacion); }

    @Override
    public String toString() {
        return String.format("%s[id=%s, nombre=%s]", getRol(), identificacion, nombre);
    }
}

// ─────────────────────────────────────────────────────────────────────────────

/**
 * Tipos de vinculación docente según el Acuerdo 004 de 2025.
 */
enum TipoVinculacion {
    /** Docente de planta con nombramiento permanente. */
    PLANTA,
    /** Docente ocasional con contrato a término definido. */
    OCASIONAL,
    /** Docente vinculado por hora cátedra. */
    HORA_CATEDRA,
    /** Docente visitante proveniente de otra institución. */
    VISITANTE,
    /** Experto o experta invitada del sector productivo. */
    EXPERTO
}
\end{lstlisting}

\begin{lstlisting}[style=javaStyle, caption={Clase \texttt{Docente} y \texttt{Estudiante}}, label={lst:docente-estudiante}]
package co.edu.udistrital.estatuto.modelo;

import java.util.Objects;

/**
 * Representa a un docente de la Universidad Distrital.
 *
 * <p>Los docentes se clasifican según su tipo de vinculación y están
 * adscritos a una Escuela. La categoría docente (Auxiliar, Asistente,
 * Asociado, Titular) aplica únicamente a los docentes de planta.</p>
 */
public class Docente extends Persona {

    private TipoVinculacion tipoVinculacion;
    private String           escuelaCodigo;
    private String           categoria;

    public Docente(String nombre, String identificacion, String correo,
                   TipoVinculacion tipo, String escuelaCodigo, String categoria) {
        super(nombre, identificacion, correo);
        this.tipoVinculacion = Objects.requireNonNull(tipo);
        this.escuelaCodigo   = escuelaCodigo;
        this.categoria       = categoria;
    }

    @Override public String getRol() { return "Docente"; }

    public TipoVinculacion getTipoVinculacion()         { return tipoVinculacion; }
    public void            setTipoVinculacion(TipoVinculacion t) { this.tipoVinculacion = t; }

    public String getEscuelaCodigo()          { return escuelaCodigo; }
    public void   setEscuelaCodigo(String ec) { this.escuelaCodigo = ec; }

    public String getCategoria()        { return categoria; }
    public void   setCategoria(String c){ this.categoria = c; }

    @Override
    public String toString() {
        return String.format("Docente[id=%s, nombre=%s, tipo=%s, escuela=%s]",
                             getIdentificacion(), getNombre(),
                             tipoVinculacion, escuelaCodigo);
    }
}

// ─────────────────────────────────────────────────────────────────────────────

/**
 * Estado académico del estudiante a lo largo de su ciclo de vida institucional.
 */
enum EstadoEstudiante {
    ASPIRANTE, ADMITIDO, ACTIVO, SUSPENDIDO, GRADUADO, RETIRADO
}

// ─────────────────────────────────────────────────────────────────────────────

/**
 * Representa a un estudiante matriculado o con historia en la Universidad.
 */
public class Estudiante extends Persona {

    private String          codigoEstudiantil;
    private String          programaCodigo;
    private int             semestre;
    private EstadoEstudiante estado;

    public Estudiante(String nombre, String identificacion, String correo,
                      String codigoEstudiantil, String programaCodigo,
                      int semestre, EstadoEstudiante estado) {
        super(nombre, identificacion, correo);
        this.codigoEstudiantil = Objects.requireNonNull(codigoEstudiantil);
        this.programaCodigo    = programaCodigo;
        this.semestre          = semestre;
        this.estado            = estado;
    }

    @Override public String getRol() { return "Estudiante"; }

    public String          getCodigoEstudiantil() { return codigoEstudiantil; }
    public String          getProgramaCodigo()    { return programaCodigo; }
    public int             getSemestre()          { return semestre; }
    public EstadoEstudiante getEstado()           { return estado; }

    public void setSemestre(int s)           { this.semestre = s; }
    public void setEstado(EstadoEstudiante e){ this.estado   = e; }
    public void setProgramaCodigo(String p)  { this.programaCodigo = p; }

    @Override
    public String toString() {
        return String.format("Estudiante[cod=%s, nombre=%s, programa=%s, semestre=%d, estado=%s]",
                             codigoEstudiantil, getNombre(), programaCodigo, semestre, estado);
    }
}
\end{lstlisting}

\begin{lstlisting}[style=javaStyle, caption={Clases \texttt{PersonalAdministrativo} y \texttt{Egresado}}, label={lst:admin-egresado}]
package co.edu.udistrital.estatuto.modelo;

/**
 * Personal administrativo de la universidad (empleados públicos
 * y trabajadores oficiales, según Art. 8, Acuerdo 004 de 2025).
 */
public class PersonalAdministrativo extends Persona {

    private String cargo;
    private String dependencia;

    public PersonalAdministrativo(String nombre, String identificacion,
                                   String correo, String cargo, String dependencia) {
        super(nombre, identificacion, correo);
        this.cargo       = cargo;
        this.dependencia = dependencia;
    }

    @Override public String getRol() { return "PersonalAdministrativo"; }

    public String getCargo()        { return cargo; }
    public void   setCargo(String c){ this.cargo = c; }

    public String getDependencia()         { return dependencia; }
    public void   setDependencia(String d) { this.dependencia = d; }
}

// ─────────────────────────────────────────────────────────────────────────────

/**
 * Egresado de la Universidad Distrital, integrante de la comunidad
 * universitaria con derechos de participación conforme al Estatuto General.
 */
public class Egresado extends Persona {

    private String programaGraduacion;
    private int    anioGraduacion;
    private String tituloObtenido;

    public Egresado(String nombre, String identificacion, String correo,
                    String programaGraduacion, int anioGraduacion, String titulo) {
        super(nombre, identificacion, correo);
        this.programaGraduacion = programaGraduacion;
        this.anioGraduacion     = anioGraduacion;
        this.tituloObtenido     = titulo;
    }

    @Override public String getRol() { return "Egresado"; }

    public String getProgramaGraduacion()         { return programaGraduacion; }
    public int    getAnioGraduacion()             { return anioGraduacion; }
    public String getTituloObtenido()             { return tituloObtenido; }
    public void   setTituloObtenido(String t)     { this.tituloObtenido = t; }
}
\end{lstlisting}

\subsection{Clases de gestión académica}
\label{subsec:gestion-academica}

El subsistema de gestión académica modela los programas, planes de estudios y espacios
académicos que conforman la oferta formativa de la universidad. Estos objetos exhiben
una cadena de composición: una \texttt{Facultad} ofrece uno o más \texttt{ProgramaAcademico},
cada uno de los cuales define un \texttt{PlanDeEstudios} que, a su vez, contiene
múltiples \texttt{EspacioAcademico}. La \texttt{Matricula} vincula al \texttt{Estudiante}
con el \texttt{ProgramaAcademico} en un periodo académico dado.

\begin{lstlisting}[style=javaStyle, caption={Clases de gestión académica}, label={lst:gestion-academica}]
package co.edu.udistrital.estatuto.modelo;

import java.util.*;

/** Nivel de formación de un programa académico. */
enum NivelPrograma {
    PREGRADO, ESPECIALIZACION, MAESTRIA, DOCTORADO
}

// ─────────────────────────────────────────────────────────────────────────────

/**
 * Programa académico ofrecido por la Universidad Distrital.
 *
 * <p>Cada programa posee un código SNIES, un nivel de formación,
 * una modalidad de oferta y un plan de estudios vigente.</p>
 */
public class ProgramaAcademico {

    private String         nombre;
    private String         codigo;
    private NivelPrograma  nivel;
    private String         modalidad;
    private String         facultadCodigo;
    private PlanDeEstudios planDeEstudios;

    public ProgramaAcademico(String nombre, String codigo,
                              NivelPrograma nivel, String modalidad,
                              String facultadCodigo) {
        this.nombre        = Objects.requireNonNull(nombre);
        this.codigo        = Objects.requireNonNull(codigo);
        this.nivel         = nivel;
        this.modalidad     = modalidad;
        this.facultadCodigo = facultadCodigo;
    }

    public String         getNombre()        { return nombre; }
    public String         getCodigo()        { return codigo; }
    public NivelPrograma  getNivel()         { return nivel; }
    public String         getModalidad()     { return modalidad; }
    public String         getFacultadCodigo(){ return facultadCodigo; }
    public PlanDeEstudios getPlanDeEstudios(){ return planDeEstudios; }

    public void setPlanDeEstudios(PlanDeEstudios p) { this.planDeEstudios = p; }
    public void setNivel(NivelPrograma n)            { this.nivel          = n; }
    public void setModalidad(String m)              { this.modalidad      = m; }

    @Override
    public String toString() {
        return String.format("Programa[%s - %s, nivel=%s]", codigo, nombre, nivel);
    }
}

// ─────────────────────────────────────────────────────────────────────────────

/**
 * Plan de estudios de un programa académico.
 * Contiene la versión vigente, los créditos totales y la lista
 * de espacios académicos que lo componen.
 */
public class PlanDeEstudios {

    private String                  version;
    private int                     creditosTotales;
    private final List<EspacioAcademico> espacios;

    public PlanDeEstudios(String version, int creditosTotales) {
        this.version        = Objects.requireNonNull(version);
        this.creditosTotales = creditosTotales;
        this.espacios       = new ArrayList<>();
    }

    public boolean agregarEspacio(EspacioAcademico ea) {
        Objects.requireNonNull(ea);
        return espacios.add(ea);
    }

    public String  getVersion()        { return version; }
    public int     getCreditosTotales(){ return creditosTotales; }

    public List<EspacioAcademico> getEspacios() {
        return Collections.unmodifiableList(espacios);
    }

    /** Suma real de créditos de los espacios registrados. */
    public int calcularCreditosAcumulados() {
        return espacios.stream().mapToInt(EspacioAcademico::getCreditos).sum();
    }
}

// ─────────────────────────────────────────────────────────────────────────────

/**
 * Espacio académico (asignatura o componente) del plan de estudios.
 */
public class EspacioAcademico {

    private String nombre;
    private String codigo;
    private int    creditos;
    private String tipo;          // p.ej. "Obligatorio", "Electivo"
    private int    horasTeoricas;
    private int    horasPracticas;

    public EspacioAcademico(String nombre, String codigo, int creditos,
                             String tipo, int horasTeoricas, int horasPracticas) {
        this.nombre        = Objects.requireNonNull(nombre);
        this.codigo        = Objects.requireNonNull(codigo);
        this.creditos      = creditos;
        this.tipo          = tipo;
        this.horasTeoricas = horasTeoricas;
        this.horasPracticas = horasPracticas;
    }

    public String getNombre()        { return nombre; }
    public String getCodigo()        { return codigo; }
    public int    getCreditos()      { return creditos; }
    public String getTipo()          { return tipo; }
    public int    getHorasTeoricas() { return horasTeoricas; }
    public int    getHorasPracticas(){ return horasPracticas; }
    public int    getHorasTotales()  { return horasTeoricas + horasPracticas; }

    @Override
    public String toString() {
        return String.format("Espacio[%s - %s, creditos=%d]", codigo, nombre, creditos);
    }
}

// ─────────────────────────────────────────────────────────────────────────────

/**
 * Matrícula que vincula a un Estudiante con un ProgramaAcademico
 * en un periodo académico determinado.
 */
public class Matricula {

    private String           estudianteId;
    private String           programaCodigo;
    private String           periodo;    // p.ej. "2025-2"
    private EstadoEstudiante estado;
    private double           valorPagado;

    public Matricula(String estudianteId, String programaCodigo,
                     String periodo, EstadoEstudiante estado) {
        this.estudianteId  = Objects.requireNonNull(estudianteId);
        this.programaCodigo = Objects.requireNonNull(programaCodigo);
        this.periodo       = Objects.requireNonNull(periodo);
        this.estado        = estado;
    }

    public String          getEstudianteId()        { return estudianteId; }
    public String          getProgramaCodigo()      { return programaCodigo; }
    public String          getPeriodo()             { return periodo; }
    public EstadoEstudiante getEstado()             { return estado; }
    public double          getValorPagado()         { return valorPagado; }

    public void setEstado(EstadoEstudiante e)       { this.estado      = e; }
    public void setValorPagado(double v)            { this.valorPagado = v; }

    @Override
    public String toString() {
        return String.format("Matricula[estudiante=%s, programa=%s, periodo=%s, estado=%s]",
                             estudianteId, programaCodigo, periodo, estado);
    }
}
\end{lstlisting}

\subsection{Clases de evaluación}
\label{subsec:evaluacion}

El sistema de evaluación formativa constituye el núcleo pedagógico del modelo.
Los Resultados de Aprendizaje (RA), articulados con el Decreto~1330 de 2019 y el
Acuerdo~02 de 2020 del CESU, se asocian a rúbricas con criterios ponderados y
niveles de dominio explícitos. Esta estructura permite que las evidencias de aprendizaje
sean registradas y calificadas de manera trazable.

\begin{lstlisting}[style=javaStyle, caption={Clases de evaluación formativa}, label={lst:evaluacion}]
package co.edu.udistrital.estatuto.modelo;

import java.time.LocalDate;
import java.util.*;

/**
 * Resultado de Aprendizaje (RA) asociado a un Programa Académico.
 *
 * <p>Formulado con verbos de acción observables (Taxonomía de Bloom)
 * y articulado con el perfil de egreso del programa.</p>
 */
public class ResultadoDeAprendizaje {

    private String descripcion;
    private String nivelDominio;   // Básico, Intermedio, Avanzado
    private String programaCodigo;

    public ResultadoDeAprendizaje(String descripcion, String nivelDominio,
                                   String programaCodigo) {
        this.descripcion   = Objects.requireNonNull(descripcion);
        this.nivelDominio  = nivelDominio;
        this.programaCodigo = programaCodigo;
    }

    public String getDescripcion()    { return descripcion; }
    public String getNivelDominio()   { return nivelDominio; }
    public String getProgramaCodigo() { return programaCodigo; }

    public void setNivelDominio(String n)  { this.nivelDominio = n; }
}

// ─────────────────────────────────────────────────────────────────────────────

/**
 * Rúbrica de evaluación con criterios y niveles de desempeño.
 */
public class Rubrica {

    private String                    nombre;
    private final List<CriterioRubrica> criterios;
    private final List<String>        niveles;  // p.ej. ["Deficiente","Aceptable","Sobresaliente"]

    public Rubrica(String nombre, List<String> niveles) {
        this.nombre    = Objects.requireNonNull(nombre);
        this.niveles   = new ArrayList<>(Objects.requireNonNull(niveles));
        this.criterios = new ArrayList<>();
    }

    public boolean agregarCriterio(CriterioRubrica c) {
        return criterios.add(Objects.requireNonNull(c));
    }

    public String              getNombre()   { return nombre; }
    public List<String>        getNiveles()  { return Collections.unmodifiableList(niveles); }
    public List<CriterioRubrica> getCriterios() {
        return Collections.unmodifiableList(criterios);
    }
}

// ─────────────────────────────────────────────────────────────────────────────

/**
 * Criterio individual dentro de una rúbrica, con descriptores por nivel.
 */
public class CriterioRubrica {

    private String              nombre;
    private double              ponderacion;   // valor entre 0 y 1
    private final Map<String, String> descriptores; // nivel -> descriptor textual

    public CriterioRubrica(String nombre, double ponderacion) {
        if (ponderacion < 0 || ponderacion > 1)
            throw new IllegalArgumentException("La ponderación debe estar entre 0 y 1.");
        this.nombre       = nombre;
        this.ponderacion  = ponderacion;
        this.descriptores = new LinkedHashMap<>();
    }

    public void agregarDescriptor(String nivel, String descriptor) {
        descriptores.put(nivel, descriptor);
    }

    public String  getNombre()      { return nombre; }
    public double  getPonderacion() { return ponderacion; }
    public Map<String, String> getDescriptores() {
        return Collections.unmodifiableMap(descriptores);
    }
}

// ─────────────────────────────────────────────────────────────────────────────

/**
 * Evidencia de aprendizaje registrada por un estudiante.
 */
public class Evidencia {

    private String      tipo;          // Tarea, Proyecto, Examen, etc.
    private String      descripcion;
    private LocalDate   fecha;
    private double      calificacion;  // 0.0 - 5.0

    public Evidencia(String tipo, String descripcion,
                     LocalDate fecha, double calificacion) {
        this.tipo          = tipo;
        this.descripcion   = descripcion;
        this.fecha         = fecha;
        setCalificacion(calificacion);
    }

    public void setCalificacion(double c) {
        if (c < 0.0 || c > 5.0)
            throw new IllegalArgumentException(
                "La calificación debe estar entre 0.0 y 5.0.");
        this.calificacion = c;
    }

    public String    getTipo()          { return tipo; }
    public String    getDescripcion()   { return descripcion; }
    public LocalDate getFecha()         { return fecha; }
    public double    getCalificacion()  { return calificacion; }
}

// ─────────────────────────────────────────────────────────────────────────────

/**
 * Propósito de Formación y Aprendizaje (PFA) en uno de los tres ámbitos
 * definidos por la Universidad Distrital (2023): ontológico, epistemológico,
 * contextual.
 */
public class PropositoFormacion {

    public enum Ambito { ONTOLOGICO, EPISTEMOLOGICO, CONTEXTUAL }

    private Ambito ambito;
    private String descripcion;
    private String programaCodigo;

    public PropositoFormacion(Ambito ambito, String descripcion, String programaCodigo) {
        this.ambito        = Objects.requireNonNull(ambito);
        this.descripcion   = Objects.requireNonNull(descripcion);
        this.programaCodigo = programaCodigo;
    }

    public Ambito getAmbito()         { return ambito; }
    public String getDescripcion()    { return descripcion; }
    public String getProgramaCodigo() { return programaCodigo; }
}
\end{lstlisting}

% ─────────────────────────────────────────────────────────────────────────────
\section{Interfaces y clases abstractas}
\label{sec:java-interfaces}
% ─────────────────────────────────────────────────────────────────────────────

\noindent
Las interfaces en \java\ cumplen un rol dual en este sistema: por un lado, contratan
el comportamiento que las clases deben exhibir; por otro, habilitan el polimorfismo
sin imponer la restricción de herencia única. Se definen tres interfaces transversales
cuya influencia se extiende a lo largo de todo el sistema.

\begin{lstlisting}[style=javaStyle, caption={Interfaces \texttt{Gestionable}, \texttt{Evaluable} y \texttt{Exportable}}, label={lst:interfaces}]
package co.edu.udistrital.estatuto.modelo;

import java.util.List;
import java.util.Optional;

/**
 * Contrato genérico para operaciones CRUD sobre entidades del dominio.
 *
 * @param <T>  tipo de la entidad gestionada
 * @param <ID> tipo del identificador (String para códigos institucionales)
 */
public interface Gestionable<T, ID> {

    /**
     * Registra una nueva entidad en el sistema.
     * @param entidad entidad a registrar
     * @return entidad registrada (puede incluir ID generado)
     */
    T crear(T entidad);

    /**
     * Recupera una entidad por su identificador.
     * @param id identificador de la entidad
     * @return {@link Optional} con la entidad, o vacío si no existe
     */
    Optional<T> buscarPorId(ID id);

    /**
     * Retorna todas las entidades del tipo gestionado.
     * @return lista inmutable de entidades
     */
    List<T> listarTodos();

    /**
     * Actualiza los datos de una entidad existente.
     * @param entidad entidad con datos actualizados
     * @return entidad actualizada
     */
    T actualizar(T entidad);

    /**
     * Elimina una entidad por su identificador.
     * @param id identificador de la entidad a eliminar
     * @return {@code true} si la eliminación fue exitosa
     */
    boolean eliminar(ID id);
}

// ─────────────────────────────────────────────────────────────────────────────

/**
 * Contrato para entidades sujetas a evaluación con rúbricas.
 */
public interface Evaluable {

    /**
     * Evalúa la entidad usando la rúbrica proporcionada.
     * @param rubrica rúbrica de evaluación
     * @return calificación resultante en escala 0.0 - 5.0
     */
    double evaluar(Rubrica rubrica);

    /**
     * Retorna las evidencias asociadas a esta entidad.
     * @return lista de evidencias registradas
     */
    List<Evidencia> getEvidencias();

    /**
     * Determina si la entidad alcanza el umbral mínimo de aprobación.
     * @param umbral calificación mínima requerida
     * @return {@code true} si supera o iguala el umbral
     */
    default boolean aprueba(double umbral) {
        return evaluar(null) >= umbral;
    }
}

// ─────────────────────────────────────────────────────────────────────────────

/**
 * Contrato para entidades que pueden exportarse como reporte.
 */
public interface Exportable {

    /**
     * Exporta la representación de la entidad en el formato especificado.
     * @param formato "JSON", "CSV", "PDF", "XML"
     * @return cadena con la representación exportada
     */
    String exportar(String formato);

    /**
     * Genera el encabezado de columnas para exportación tabular.
     * @return arreglo con los nombres de las columnas
     */
    String[] getEncabezados();
}
\end{lstlisting}

% ─────────────────────────────────────────────────────────────────────────────
\section{Herencia, polimorfismo y encapsulamiento}
\label{sec:java-poo-pilares}
% ─────────────────────────────────────────────────────────────────────────────

\noindent
Los tres pilares de la \poo\ se manifiestan con nitidez en el diseño del sistema.
El polimorfismo emerge de manera natural cuando el código cliente opera sobre
referencias de tipo \texttt{UnidadAcademica} sin conocer el tipo concreto de la
instancia; el encapsulamiento protege la integridad del estado de los objetos
mediante validaciones en los mutadores; la herencia estructura la jerarquía de
entidades evitando la duplicación de código.

\begin{lstlisting}[style=javaStyle, caption={Ejemplo de polimorfismo con \texttt{UnidadAcademica}}, label={lst:polimorfismo}]
package co.edu.udistrital.estatuto.util;

import co.edu.udistrital.estatuto.modelo.*;
import java.util.List;

/**
 * Demostración de uso polimórfico de UnidadAcademica.
 * El método imprimirInfo() opera sobre cualquier subclase
 * sin requerir instanceof ni castings explícitos.
 */
public class DemoPolimorfismo {

    /**
     * Imprime la información básica de cada unidad académica.
     * El método getTipo() se resuelve en tiempo de ejecución
     * (despacho dinámico / virtual dispatch).
     *
     * @param unidades lista heterogénea de unidades académicas
     */
    public static void imprimirResumen(List<UnidadAcademica> unidades) {
        for (UnidadAcademica ua : unidades) {
            // getTipo() retorna "Facultad", "Escuela", "Centro" o "Instituto"
            // según el tipo concreto, sin que este método lo conozca.
            System.out.printf("%-12s | %-8s | %s%n",
                              ua.getTipo(), ua.getCodigo(), ua.getNombre());
        }
    }

    /** Ejemplo de encapsulamiento con validación en setter. */
    public static void demoEncapsulamiento() {
        Facultad f = new Facultad("Ingeniería", "FI",
                                  java.time.LocalDate.of(1950, 1, 1),
                                  "Dir. General", "Dr. Ejemplo");
        try {
            f.setCodigo("");            // Lanza IllegalArgumentException
        } catch (IllegalArgumentException ex) {
            System.out.println("Validación activa: " + ex.getMessage());
        }
        f.setCodigo("FI-001");          // Código válido: normalizado a mayúsculas
        System.out.println("Código actualizado: " + f.getCodigo());
    }

    public static void main(String[] args) {
        List<UnidadAcademica> unidades = List.of(
            new Facultad("Ingeniería",  "FI",  java.time.LocalDate.of(1950,1,1), "D1","D2"),
            new Escuela("Ing. Sistemas","EIS", java.time.LocalDate.of(1970,1,1), "D3"),
            new Centro("CIDC",          "CID", java.time.LocalDate.of(1995,1,1), "D4")
        );
        imprimirResumen(unidades);
        demoEncapsulamiento();
    }
}
\end{lstlisting}

\noindent
La salida del método \texttt{imprimirResumen} evidencia el polimorfismo en acción:
la misma instrucción \texttt{ua.getTipo()} retorna cadenas distintas
---\texttt{"Facultad"}, \texttt{"Escuela"}, \texttt{"Centro"}--- según el tipo
concreto de cada objeto. No hay condicionales; el compilador y la JVM resuelven el
método correcto en tiempo de ejecución mediante el mecanismo de despacho dinámico.
Esto es, precisamente, lo que distingue a la orientación a objetos del paradigma
procedural: el \textit{qué hacer} se delega al objeto mismo.

% ─────────────────────────────────────────────────────────────────────────────
\section{Patrones de diseño}
\label{sec:java-patrones}
% ─────────────────────────────────────────────────────────────────────────────

\noindent
Los patrones de diseño ---catálogo sistematizado por \textcite{Gamma1994}--- ofrecen
soluciones recurrentes a problemas recurrentes de diseño orientado a objetos. Se
implementan tres patrones que responden a necesidades concretas del dominio
estatutario: Factory Method, Strategy y Observer.

\subsection{Factory Method: creación de unidades académicas}

El patrón Factory Method desacopla la creación de objetos de su tipo concreto.
En este sistema, la \texttt{UnidadAcademicaFactory} centraliza la instanciación
de Facultades, Escuelas, Centros e Institutos, evitando que el código cliente
use constructores directamente y facilitando la extensión del sistema con nuevos
tipos de unidad sin modificar el código existente.\footnote{Este comportamiento
corresponde al principio Abierto/Cerrado (OCP) de SOLID: el sistema está abierto
a la extensión (nuevos tipos de unidad) pero cerrado a la modificación del código cliente.}

\begin{lstlisting}[style=javaStyle, caption={Factory Method para \texttt{UnidadAcademica}}, label={lst:factory}]
package co.edu.udistrital.estatuto.patron;

import co.edu.udistrital.estatuto.modelo.*;
import java.time.LocalDate;

/**
 * Factory Method para la creación de unidades académicas.
 *
 * <p>Centraliza la instanciación y permite incorporar lógica de
 * inicialización compartida (p. ej., auditoría, logging) sin duplicar
 * código en cada punto de creación.</p>
 */
public class UnidadAcademicaFactory {

    public enum TipoUnidad { FACULTAD, ESCUELA, CENTRO, INSTITUTO }

    /**
     * Crea una unidad académica del tipo especificado.
     *
     * @param tipo          tipo de unidad a crear
     * @param nombre        nombre oficial
     * @param codigo        código institucional
     * @param fechaCreacion fecha de creación
     * @param director      nombre del director
     * @return instancia concreta de {@link UnidadAcademica}
     * @throws IllegalArgumentException si el tipo no es reconocido
     */
    public static UnidadAcademica crear(TipoUnidad tipo,
                                        String nombre, String codigo,
                                        LocalDate fechaCreacion, String director) {
        return switch (tipo) {
            case FACULTAD  -> new Facultad(nombre, codigo, fechaCreacion,
                                           director, director);
            case ESCUELA   -> new Escuela(nombre, codigo, fechaCreacion, director);
            case CENTRO    -> new Centro(nombre, codigo, fechaCreacion, director);
            case INSTITUTO -> new Instituto(nombre, codigo, fechaCreacion, director);
        };
    }

    // Ejemplo de uso
    public static void main(String[] args) {
        UnidadAcademica fi = UnidadAcademicaFactory.crear(
            TipoUnidad.FACULTAD, "Facultad de Ingeniería",
            "FI", LocalDate.of(1950, 1, 1), "Decano Ejemplo");

        UnidadAcademica eis = UnidadAcademicaFactory.crear(
            TipoUnidad.ESCUELA, "Escuela de Ingeniería de Sistemas",
            "EIS", LocalDate.of(1970, 6, 15), "Director Ejemplo");

        System.out.println(fi);
        System.out.println(eis);
    }
}
\end{lstlisting}

\subsection{Strategy: algoritmos de evaluación intercambiables}

El patrón Strategy encapsula algoritmos alternativos detrás de una interfaz común,
permitiendo seleccionar el algoritmo en tiempo de ejecución. En el contexto del
sistema de evaluación, distintos programas académicos pueden requerir estrategias
de cálculo de nota final disímiles: promedio aritmético, promedio ponderado por
créditos, o escala de suficiencia/proficiencia.

\begin{lstlisting}[style=javaStyle, caption={Strategy para algoritmos de evaluación}, label={lst:strategy}]
package co.edu.udistrital.estatuto.patron;

import co.edu.udistrital.estatuto.modelo.Evidencia;
import java.util.List;

/**
 * Interfaz Strategy para el cálculo de nota final a partir
 * de un conjunto de evidencias de aprendizaje.
 */
public interface EvaluacionStrategy {

    /**
     * Calcula la nota final con base en las evidencias proporcionadas.
     *
     * @param evidencias lista de evidencias con sus calificaciones individuales
     * @return nota final en escala 0.0 - 5.0
     */
    double calcularNotaFinal(List<Evidencia> evidencias);

    /** Nombre descriptivo de la estrategia (para reportes). */
    String getNombreEstrategia();
}

// ─────────────────────────────────────────────────────────────────────────────

/** Estrategia 1: promedio aritmético simple. */
class PromedioAritmeticoStrategy implements EvaluacionStrategy {

    @Override
    public double calcularNotaFinal(List<Evidencia> evidencias) {
        if (evidencias == null || evidencias.isEmpty()) return 0.0;
        return evidencias.stream()
                         .mapToDouble(Evidencia::getCalificacion)
                         .average()
                         .orElse(0.0);
    }

    @Override
    public String getNombreEstrategia() { return "Promedio Aritmético"; }
}

// ─────────────────────────────────────────────────────────────────────────────

/** Estrategia 2: nota mínima de tres (nota más baja de todas las evidencias). */
class NotaMinimaStrategy implements EvaluacionStrategy {

    @Override
    public double calcularNotaFinal(List<Evidencia> evidencias) {
        if (evidencias == null || evidencias.isEmpty()) return 0.0;
        return evidencias.stream()
                         .mapToDouble(Evidencia::getCalificacion)
                         .min()
                         .orElse(0.0);
    }

    @Override
    public String getNombreEstrategia() { return "Nota Mínima"; }
}

// ─────────────────────────────────────────────────────────────────────────────

/**
 * Contexto que utiliza la estrategia de evaluación.
 * Permite cambiar la estrategia en tiempo de ejecución.
 */
class CalculadorDeNota {

    private EvaluacionStrategy estrategia;

    public CalculadorDeNota(EvaluacionStrategy estrategia) {
        this.estrategia = estrategia;
    }

    /** Cambia la estrategia sin modificar el contexto. */
    public void setEstrategia(EvaluacionStrategy nueva) {
        this.estrategia = nueva;
    }

    public double calcular(List<Evidencia> evidencias) {
        return estrategia.calcularNotaFinal(evidencias);
    }

    public void imprimirResultado(List<Evidencia> evidencias) {
        System.out.printf("Estrategia: %-25s | Nota: %.2f%n",
                          estrategia.getNombreEstrategia(),
                          calcular(evidencias));
    }
}
\end{lstlisting}

\subsection{Observer: notificaciones del Consejo}

El patrón Observer implementa una dependencia uno-a-muchos entre objetos: cuando
el estado del sujeto cambia, todos sus observadores son notificados automáticamente.
En el sistema estatutario, el Consejo de Facultad actúa como sujeto observable;
los interesados en sus decisiones ---sistemas de notificación, módulos de auditoría---
se suscriben como observadores.

\begin{lstlisting}[style=javaStyle, caption={Observer para notificaciones del Consejo}, label={lst:observer}]
package co.edu.udistrital.estatuto.patron;

import java.util.*;

/**
 * Interface Observer: observador de eventos del Consejo.
 */
public interface ConsejoObserver {
    /**
     * Notifica al observador sobre un evento del Consejo.
     * @param evento     descripción del evento ocurrido
     * @param consejo    nombre del consejo que generó el evento
     */
    void actualizar(String evento, String consejo);
}

// ─────────────────────────────────────────────────────────────────────────────

/**
 * Consejo observable: notifica a todos los observadores suscritos
 * cuando se registra un nuevo acuerdo o decisión.
 */
class ConsejoObservable {

    private final String              nombre;
    private final List<ConsejoObserver> observadores;
    private String                    ultimoAcuerdo;

    public ConsejoObservable(String nombre) {
        this.nombre       = nombre;
        this.observadores = new ArrayList<>();
    }

    /** Suscribe un nuevo observador. */
    public void suscribir(ConsejoObserver obs) {
        observadores.add(Objects.requireNonNull(obs));
    }

    /** Cancela la suscripción de un observador. */
    public void desuscribir(ConsejoObserver obs) {
        observadores.remove(obs);
    }

    /**
     * Registra un nuevo acuerdo y notifica a todos los observadores.
     * @param acuerdo texto del acuerdo aprobado
     */
    public void registrarAcuerdo(String acuerdo) {
        this.ultimoAcuerdo = acuerdo;
        notificar("ACUERDO_APROBADO: " + acuerdo);
    }

    private void notificar(String evento) {
        for (ConsejoObserver obs : observadores) {
            obs.actualizar(evento, nombre);
        }
    }

    public String getUltimoAcuerdo() { return ultimoAcuerdo; }
}

// ─────────────────────────────────────────────────────────────────────────────

/** Observador concreto: módulo de auditoría. */
class AuditoriaObserver implements ConsejoObserver {
    private final List<String> bitacora = new ArrayList<>();

    @Override
    public void actualizar(String evento, String consejo) {
        String registro = String.format("[AUDITORIA] %s - Consejo: %s", evento, consejo);
        bitacora.add(registro);
        System.out.println(registro);
    }

    public List<String> getBitacora() {
        return Collections.unmodifiableList(bitacora);
    }
}
\end{lstlisting}

% ─────────────────────────────────────────────────────────────────────────────
\section{Manejo de excepciones}
\label{sec:java-excepciones}
% ─────────────────────────────────────────────────────────────────────────────

\noindent
El manejo estructurado de excepciones constituye un aspecto crítico de los sistemas
de información institucional, donde la corrupción silenciosa de datos puede tener
consecuencias administrativas y legales. Se define una jerarquía de excepciones
encabezada por \texttt{EstatutoException}, clase verificada (\textit{checked}) que
obliga a los clientes del sistema a tratar explícitamente los errores del dominio.

\begin{lstlisting}[style=javaStyle, caption={Jerarquía de excepciones del sistema}, label={lst:excepciones}]
package co.edu.udistrital.estatuto.excepcion;

/**
 * Excepción raíz del sistema de información estatutario.
 * Todas las excepciones de negocio heredan de esta clase.
 */
public class EstatutoException extends Exception {

    private final String codigoError;

    public EstatutoException(String mensaje, String codigoError) {
        super(mensaje);
        this.codigoError = codigoError;
    }

    public EstatutoException(String mensaje, String codigoError, Throwable causa) {
        super(mensaje, causa);
        this.codigoError = codigoError;
    }

    public String getCodigoError() { return codigoError; }
}

// ─────────────────────────────────────────────────────────────────────────────

/**
 * Excepción lanzada cuando una operación de matrícula es inválida.
 * Ejemplos: matrícula en programa cerrado, doble matrícula, saldo pendiente.
 */
class MatriculaException extends EstatutoException {

    public MatriculaException(String mensaje) {
        super(mensaje, "MAT-001");
    }

    public MatriculaException(String mensaje, String codigo) {
        super(mensaje, codigo);
    }
}

// ─────────────────────────────────────────────────────────────────────────────

/**
 * Excepción lanzada cuando una operación de evaluación es inválida.
 */
class EvaluacionException extends EstatutoException {

    public EvaluacionException(String mensaje) {
        super(mensaje, "EVA-001");
    }
}

// ─────────────────────────────────────────────────────────────────────────────

/**
 * Excepción lanzada cuando una entidad buscada no existe en el sistema.
 */
class EntidadNoEncontradaException extends EstatutoException {

    private final String tipo;
    private final String id;

    public EntidadNoEncontradaException(String tipo, String id) {
        super(String.format("No se encontró %s con identificador '%s'.", tipo, id),
              "ENT-404");
        this.tipo = tipo;
        this.id   = id;
    }

    public String getTipoEntidad() { return tipo; }
    public String getIdBuscado()   { return id; }
}
\end{lstlisting}

% ─────────────────────────────────────────────────────────────────────────────
\section{Interfaz gráfica con JavaFX}
\label{sec:java-javafx}
% ─────────────────────────────────────────────────────────────────────────────

\noindent
La interfaz gráfica del sistema adopta una estructura de tres niveles: una barra
de menú (\texttt{MenuBar}) para las operaciones globales, un panel de pestañas
(\texttt{TabPane}) para gestionar las entidades del dominio (Facultades, Docentes,
Estudiantes, Programas) y, dentro de cada pestaña, una tabla de resultados
(\texttt{TableView}) con botones CRUD. Esta estructura refleja el patrón MVC:
las vistas delegan toda la lógica en el controlador y no acceden directamente
al modelo de datos.

\begin{lstlisting}[style=javaStyle, caption={Vista principal con JavaFX}, label={lst:javafx-vista}]
package co.edu.udistrital.estatuto.vista;

import co.edu.udistrital.estatuto.controlador.FacultadControlador;
import co.edu.udistrital.estatuto.modelo.Facultad;
import javafx.application.Application;
import javafx.beans.property.SimpleStringProperty;
import javafx.collections.FXCollections;
import javafx.collections.ObservableList;
import javafx.scene.Scene;
import javafx.scene.control.*;
import javafx.scene.layout.*;
import javafx.stage.Stage;

/**
 * Vista principal de la aplicación JavaFX.
 *
 * <p>Organiza la interfaz en un TabPane con pestañas para cada
 * entidad principal del dominio. La pestaña de Facultades se
 * implementa de forma completa; las demás siguen el mismo patrón.</p>
 */
public class VistaPrincipal extends Application {

    private final FacultadControlador ctrlFacultad = new FacultadControlador();

    @Override
    public void start(Stage stage) {
        BorderPane raiz = new BorderPane();

        // ── Barra de menú ───────────────────────────────────────────────────
        MenuBar menuBar = crearMenuBar(stage);
        raiz.setTop(menuBar);

        // ── Panel de pestañas ───────────────────────────────────────────────
        TabPane tabPane = new TabPane();
        tabPane.getTabs().addAll(
            crearPestaniaFacultades(),
            crearPestaniaPlaceholder("Docentes"),
            crearPestaniaPlaceholder("Estudiantes"),
            crearPestaniaPlaceholder("Programas")
        );
        raiz.setCenter(tabPane);

        Scene escena = new Scene(raiz, 1024, 700);
        stage.setTitle("Sistema de Información Estatutario — Universidad Distrital");
        stage.setScene(escena);
        stage.show();
    }

    // ── Menú principal ──────────────────────────────────────────────────────

    private MenuBar crearMenuBar(Stage stage) {
        Menu menuArchivo = new Menu("Archivo");
        MenuItem miSalir = new MenuItem("Salir");
        miSalir.setOnAction(e -> stage.close());
        menuArchivo.getItems().add(miSalir);

        Menu menuAyuda = new Menu("Ayuda");
        MenuItem miAcerca = new MenuItem("Acerca de...");
        miAcerca.setOnAction(e -> mostrarDialogoAcerca());
        menuAyuda.getItems().add(miAcerca);

        return new MenuBar(menuArchivo, menuAyuda);
    }

    // ── Pestaña de Facultades ───────────────────────────────────────────────

    private Tab crearPestaniaFacultades() {
        Tab tab = new Tab("Facultades");
        tab.setClosable(false);

        // Tabla de Facultades
        TableView<Facultad> tabla = new TableView<>();
        ObservableList<Facultad> datos = FXCollections.observableArrayList(
            ctrlFacultad.listarTodos()
        );
        tabla.setItems(datos);

        TableColumn<Facultad, String> colCodigo = new TableColumn<>("Código");
        colCodigo.setCellValueFactory(
            c -> new SimpleStringProperty(c.getValue().getCodigo()));

        TableColumn<Facultad, String> colNombre = new TableColumn<>("Nombre");
        colNombre.setCellValueFactory(
            c -> new SimpleStringProperty(c.getValue().getNombre()));

        TableColumn<Facultad, String> colDecano = new TableColumn<>("Decano/a");
        colDecano.setCellValueFactory(
            c -> new SimpleStringProperty(c.getValue().getDecano()));

        tabla.getColumns().addAll(colCodigo, colNombre, colDecano);

        // Botones CRUD
        Button btnNuevo    = new Button("Nueva Facultad");
        Button btnEditar   = new Button("Editar");
        Button btnEliminar = new Button("Eliminar");

        btnNuevo.setOnAction(e -> mostrarFormularioFacultad(null, datos));
        btnEditar.setOnAction(e -> {
            Facultad sel = tabla.getSelectionModel().getSelectedItem();
            if (sel != null) mostrarFormularioFacultad(sel, datos);
        });
        btnEliminar.setOnAction(e -> {
            Facultad sel = tabla.getSelectionModel().getSelectedItem();
            if (sel != null && ctrlFacultad.eliminar(sel.getCodigo())) {
                datos.remove(sel);
            }
        });

        HBox botones = new HBox(8, btnNuevo, btnEditar, btnEliminar);
        VBox contenido = new VBox(8, tabla, botones);
        contenido.setPadding(new javafx.geometry.Insets(10));
        tab.setContent(contenido);
        return tab;
    }

    private Tab crearPestaniaPlaceholder(String nombre) {
        Tab tab = new Tab(nombre);
        tab.setClosable(false);
        tab.setContent(new Label("Sección " + nombre + " — en construcción"));
        return tab;
    }

    private void mostrarFormularioFacultad(Facultad f,
                                            ObservableList<Facultad> datos) {
        Dialog<Facultad> dialog = new Dialog<>();
        dialog.setTitle(f == null ? "Nueva Facultad" : "Editar Facultad");
        TextField tfCodigo  = new TextField(f != null ? f.getCodigo()  : "");
        TextField tfNombre  = new TextField(f != null ? f.getNombre()  : "");
        TextField tfDecano  = new TextField(f != null ? f.getDecano()  : "");
        GridPane grid = new GridPane();
        grid.setHgap(8); grid.setVgap(8);
        grid.addRow(0, new Label("Código:"), tfCodigo);
        grid.addRow(1, new Label("Nombre:"), tfNombre);
        grid.addRow(2, new Label("Decano/a:"), tfDecano);
        dialog.getDialogPane().setContent(grid);
        dialog.getDialogPane().getButtonTypes().addAll(
            ButtonType.OK, ButtonType.CANCEL);
        dialog.setResultConverter(bt -> {
            if (bt == ButtonType.OK) {
                // Aquí se invoca el controlador, no el modelo directamente
                return ctrlFacultad.crearOActualizar(
                    tfCodigo.getText(), tfNombre.getText(), tfDecano.getText());
            }
            return null;
        });
        dialog.showAndWait().ifPresent(nueva -> {
            if (!datos.contains(nueva)) datos.add(nueva);
        });
    }

    private void mostrarDialogoAcerca() {
        Alert alert = new Alert(Alert.AlertType.INFORMATION);
        alert.setTitle("Acerca de");
        alert.setHeaderText("Sistema de Información Estatutario");
        alert.setContentText("Desarrollado como caso de estudio del libro POO.\n"
                           + "Universidad Distrital Francisco José de Caldas, 2026.");
        alert.showAndWait();
    }
}
\end{lstlisting}

% ─────────────────────────────────────────────────────────────────────────────
\section{Persistencia de datos}
\label{sec:java-persistencia}
% ─────────────────────────────────────────────────────────────────────────────

\noindent
La persistencia se implementa mediante el patrón DAO (\textit{Data Access Object}),
que desacopla la lógica de acceso a datos de las clases de dominio. Los objetos se
serializan a formato JSON utilizando la biblioteca Gson, cuya simplicidad la hace
apropiada para el alcance de este proyecto de enseñanza.\footnote{En sistemas de
producción se preferiría un ORM como Hibernate o un cliente JPA. Para los propósitos
pedagógicos de este libro, la serialización JSON con Gson ofrece la virtud de la
transparencia: el estudiante puede abrir el archivo resultante y verificar su contenido
sin herramientas adicionales.}

\begin{lstlisting}[style=javaStyle, caption={Interfaz \texttt{GenericDAO} y clase \texttt{JsonDAO}}, label={lst:dao}]
package co.edu.udistrital.estatuto.persistencia;

import com.google.gson.*;
import com.google.gson.reflect.TypeToken;
import java.io.*;
import java.lang.reflect.Type;
import java.nio.file.*;
import java.util.*;

/**
 * Interfaz genérica para operaciones de acceso a datos.
 *
 * @param <T>  tipo de la entidad
 * @param <ID> tipo del identificador
 */
public interface GenericDAO<T, ID> {
    void       guardar(T entidad)    throws IOException;
    Optional<T> buscarPorId(ID id)  throws IOException;
    List<T>    buscarTodos()         throws IOException;
    void       actualizar(T entidad) throws IOException;
    boolean    eliminar(ID id)       throws IOException;
}

// ─────────────────────────────────────────────────────────────────────────────

/**
 * Implementación de DAO con serialización JSON (Gson).
 * Persiste listas de entidades en archivos .json en el directorio de datos.
 *
 * @param <T> tipo de la entidad a persistir
 */
public class JsonDAO<T> {

    private final Gson   gson;
    private final Path   archivo;
    private final Type   tipoLista;

    /**
     * @param rutaArchivo ruta del archivo JSON de datos
     * @param tipoLista   tipo de la lista, p.ej. {@code new TypeToken<List<Facultad>>(){}.getType()}
     */
    public JsonDAO(String rutaArchivo, Type tipoLista) {
        this.gson      = new GsonBuilder()
                              .registerTypeAdapter(
                                  java.time.LocalDate.class,
                                  (JsonSerializer<java.time.LocalDate>)
                                      (src, t, ctx) -> new JsonPrimitive(src.toString()))
                              .setPrettyPrinting()
                              .create();
        this.archivo   = Path.of(rutaArchivo);
        this.tipoLista = tipoLista;
    }

    /** Lee la lista completa del archivo JSON. */
    @SuppressWarnings("unchecked")
    public <T2> List<T2> leerTodos() throws IOException {
        if (!Files.exists(archivo)) return new ArrayList<>();
        try (Reader r = Files.newBufferedReader(archivo)) {
            List<T2> lista = gson.fromJson(r, tipoLista);
            return lista != null ? lista : new ArrayList<>();
        }
    }

    /** Escribe la lista completa en el archivo JSON. */
    public <T2> void guardarTodos(List<T2> lista) throws IOException {
        Files.createDirectories(archivo.getParent());
        try (Writer w = Files.newBufferedWriter(archivo)) {
            gson.toJson(lista, tipoLista, w);
        }
    }
}
\end{lstlisting}

% ─────────────────────────────────────────────────────────────────────────────
%\section{\java\ versus \rust: análisis comparativo de implementación}

\section[Java versus Rust: análisis comparativo]{\java\ versus \rust: análisis comparativo de implementación}
\label{sec:java-vs-rust}
% ─────────────────────────────────────────────────────────────────────────────

\noindent
La implementación paralela del mismo dominio en dos lenguajes tan disímiles como
\java\ y \rust\ suscita reflexiones que no emergen cuando se estudia un solo lenguaje.
La tabla~\ref{tab:java-rust-comparacion} sintetiza las diferencias más relevantes
desde la perspectiva de la implementación concreta del sistema estatutario.

%{\shorthandoff{<>"'`}
	
\begingroup
\small
\begin{longtable}{@{} L{2.8cm} L{5.2cm} L{5.5cm} @{}}
	% --- Título y encabezado de la primera página ---
	\caption{Comparación de implementación \java\ vs.\ \rust\ en el sistema estatutario.} \label{tab:java-rust-comparacion} \\
	\toprule
	\textbf{Dimensión} & \textbf{\java} & \textbf{\rust} \\
	\midrule
	\endfirsthead
	
	% --- Encabezado para las páginas siguientes ---
	\multicolumn{3}{c}{{\bfseries \tablename\ \thetable{} -- Continuación}} \\
	\toprule
	\textbf{Dimensión} & \textbf{\java} & \textbf{\rust} \\
	\midrule
	\endhead
	
	% --- Pie de tabla para páginas intermedias ---
	\midrule
	\multicolumn{3}{r}{{Continúa en la siguiente página...}} \\
	\bottomrule
	\endfoot
	
	% --- Pie de tabla final ---
	\bottomrule
	\endlastfoot
	
	% --- Datos de la tabla ---
	Herencia de clases &
	Herencia simple de clases + implementación múltiple de interfaces. \texttt{Facultad extends UnidadAcademica} &
	No existe herencia de structs. Se usan traits + composición. Un struct \texttt{Facultad} implementa el trait \texttt{EntidadAcademica} \\
	\midrule
	
	Manejo de \texttt{null} &
	\texttt{null} es un valor válido para cualquier referencia; riesgo de \texttt{NullPointerException}. Se mitiga con \texttt{Objects.requireNonNull} &
	No existe \texttt{null}. Se usa \verb|Option<T>|: \verb|Some(valor)| o \verb|None|. El compilador obliga a tratar ambos casos \\
	\midrule
	
	Excepciones &
	Excepciones checked y unchecked. La jerarquía \texttt{EstatutoException} impone tratamiento explícito &
	No hay excepciones. Se usa \verb|Result<T, E>|: \verb|Ok(valor)| o \verb|Err(error)|. Propagación con operador \verb|?| \\
	\midrule
	
	Concurrencia &
	\texttt{synchronized}, \texttt{ReentrantLock}, \texttt{ExecutorService}. Errores de carrera detectados en ejecución &
	\verb|Arc<Mutex<T>>|. El compilador garantiza que \verb|T| implemente \verb|Send + Sync|. Errores de carrera detectados en compilación \\
\end{longtable}
\endgroup
	
%} % ← \shorthandon implícito al cerrar el grupo

\noindent
Un aspecto particularmente ilustrativo es el tratamiento del valor ausente.
En \java, la búsqueda de un \texttt{Estudiante} por código puede retornar \texttt{null},
con el consecuente riesgo de \texttt{NullPointerException} si el llamador omite la
verificación. En \rust, la función retorna \texttt{Option<\&Estudiante>} y el compilador
obliga al llamador a tratar tanto el caso \texttt{Some(\&e)} como el caso \texttt{None}
mediante un \texttt{match} exhaustivo. El error no puede ocurrir en ejecución porque
no puede compilar sin tratamiento explícito. Esta diferencia no es menor:
habida cuenta de que los sistemas de información universitaria operan con datos críticos
de matrícula y evaluación, la garantía en tiempo de compilación resulta especialmente
valiosa.

% ─────────────────────────────────────────────────────────────────────────────
\section*{Sprints 6 y 7: implementación en \java}
\label{sec:sprints-6-7}
\addcontentsline{toc}{section}{Sprints 6 y 7: implementación en \java}
% ─────────────────────────────────────────────────────────────────────────────

\begin{tcolorbox}[sprintBox, title={Sprint 6 — Clases del dominio y pruebas unitarias}]
\textbf{Duración:} 2 semanas\\[4pt]
\textbf{Objetivo:} Implementar las clases del dominio (modelo) con cobertura de pruebas unitarias $\geq 80\%$.

\vspace{6pt}
\begin{tabularx}{\linewidth}{@{}L{4cm}L{4cm}L{4.5cm}@{}}
\toprule
\textbf{Entregable técnico} & \textbf{Característica del programa} & \textbf{Criterio de la rúbrica} \\
\midrule
Clases del paquete \texttt{modelo} compiladas y documentadas con Javadoc &
  Clases de entidad organizacional, actores y gestión académica &
  Código compila sin errores; Javadoc presente en todas las clases públicas \\
\addlinespace
Suite de pruebas JUnit~5 para \texttt{UnidadAcademica}, \texttt{Persona} y subclases &
  Herencia, polimorfismo, encapsulamiento y validaciones &
  Cobertura $\geq 80\%$ medida con JaCoCo; todos los tests pasan \\
\addlinespace
Demostración de polimorfismo con lista de \texttt{UnidadAcademica} &
  Despacho dinámico y patrón Factory Method &
  El método \texttt{getTipo()} retorna el valor correcto para cada subclase \\
\addlinespace
Jerarquía de excepciones \texttt{EstatutoException} completa &
  Manejo estructurado de errores &
  Las excepciones llevan código de error; los mensajes son descriptivos \\
\bottomrule
\end{tabularx}

\vspace{8pt}
\textbf{Definición de terminado:} El proyecto Maven compila con \texttt{mvn clean test};
todos los tests pasan; el reporte JaCoCo muestra cobertura $\geq 80\%$;
el Javadoc se genera sin advertencias con \texttt{mvn javadoc:javadoc}.
\end{tcolorbox}

\vspace{10pt}

\begin{tcolorbox}[actividadIA, title={Actividad IA — Sprint 6: Generación de pruebas unitarias}]
\textbf{Contexto:} Se propone explorar el uso de asistentes generativos como auxiliares
en la escritura de pruebas unitarias, tarea que los estudiantes suelen postergar por
considerarla tediosa pero que resulta decisiva para la calidad del software.

\vspace{6pt}
\textbf{Tarea concreta:}
\begin{enumerate}
  \item Solicitar a un asistente generativo: \textit{``Genera una clase de prueba JUnit~5
  para la clase \texttt{Escuela} del paquete \texttt{co.edu.udistrital.estatuto.modelo}.
  Incluye pruebas para: agregar docente duplicado (debe retornar \texttt{false}),
  agregar docente nulo (debe lanzar \texttt{NullPointerException}),
  y verificar que \texttt{getTipo()} retorne exactamente la cadena \textit{`Escuela'}'.}
  \item Integrar el código generado al proyecto y ejecutarlo con \texttt{mvn test}.
  \item Identificar al menos dos casos de prueba que el asistente omitió.
  \item Escribir manualmente esos casos y documentar el proceso en el informe.
\end{enumerate}

\vspace{6pt}
\textbf{Reflexión:} ¿El asistente generativo prueba el comportamiento esperado
o solamente el comportamiento que el código ya exhibe? ¿Cuál es la diferencia
entre ambos enfoques?

\vspace{6pt}
\textbf{Advertencia epistémica:} Un asistente generativo no conoce el \acuerdo\
ni las reglas de negocio específicas de la \ud. Las pruebas que genera son
genéricas. La responsabilidad de identificar los casos límite que importan
---como la restricción de que una Escuela no admita docentes duplicados--- recae
exclusivamente en el programador.
\end{tcolorbox}

\vspace{10pt}

\begin{tcolorbox}[sprintBox, title={Sprint 7 — Interfaz gráfica, persistencia y patrones de diseño}]
\textbf{Duración:} 2 semanas\\[4pt]
\textbf{Objetivo:} Integrar la capa de vista (JavaFX), la persistencia JSON y al menos
dos patrones de diseño documentados.

\vspace{6pt}
\begin{tabularx}{\linewidth}{@{}L{4cm}L{4cm}L{4.5cm}@{}}
\toprule
\textbf{Entregable técnico} & \textbf{Característica del programa} & \textbf{Criterio de la rúbrica} \\
\midrule
Aplicación JavaFX con ventana principal funcional &
  TabPane con pestaña de Facultades operativa (CRUD completo) &
  La aplicación inicia sin errores; las operaciones CRUD persisten al reiniciar \\
\addlinespace
Persistencia JSON con \texttt{JsonDAO} para Facultades &
  Serialización/deserialización con Gson; archivos \texttt{.json} en directorio \texttt{datos/} &
  Los objetos se recuperan correctamente tras reiniciar la aplicación \\
\addlinespace
Implementación documentada de Factory Method &
  Creación polimórfica de unidades académicas &
  El código usa \texttt{UnidadAcademicaFactory.crear()} en lugar de constructores directos \\
\addlinespace
Implementación documentada de Observer &
  Notificación automática de cambios en el Consejo &
  Al menos dos observadores concretos suscritos; la bitácora registra todos los eventos \\
\addlinespace
Informe técnico de 3 páginas &
  Justificación de decisiones de diseño &
  El informe explica por qué se eligió cada patrón en relación al dominio estatutario \\
\bottomrule
\end{tabularx}

\vspace{8pt}
\textbf{Definición de terminado:} La aplicación empaquetada como JAR ejecutable arranca,
muestra la lista de Facultades guardadas previamente y permite agregar/editar/eliminar.
Los patrones de diseño están implementados y su uso se demuestra en la aplicación.
\end{tcolorbox}

\vspace{10pt}

\begin{tcolorbox}[actividadIA, title={Actividad IA — Sprint 7: Patrones de diseño con asistencia generativa}]
\textbf{Tarea:}
\begin{enumerate}
  \item Solicitar a un asistente generativo la implementación del patrón Strategy
  para el cálculo de nota final. Especificar: \textit{``Implementa el patrón Strategy
  en Java 17 para calcular la nota final de un estudiante. Las estrategias deben ser:
  promedio aritmético, promedio ponderado (cada evidencia tiene un peso de 0 a 1),
  y nota máxima. Usa la clase \texttt{Evidencia} ya definida en el dominio.''}
  \item Comparar la solución generada con la implementada en la sección~\ref{sec:java-patrones}.
  \item Identificar diferencias en nomenclatura, estructura y cobertura de casos.
  \item Evaluar: ¿el asistente añadió estrategias que usted no había considerado?
  ¿eliminó alguna restricción del dominio?
\end{enumerate}

\vspace{6pt}
\textbf{Reflexión crítica:} El valor de un asistente generativo en la implementación de
patrones de diseño reside en la velocidad de andamiaje (\textit{scaffolding}), no en
la comprensión del dominio. El patrón puede ser sintácticamente correcto pero
semánticamente inadecuado si el asistente desconoce las reglas de negocio.
La revisión crítica de la solución generada es, por ello, una competencia irrenunciable.
\end{tcolorbox}
